\chapter{Positionnement et sentiment}\label{ch:s1:positioning-and-sentiment}

Chaque vendredi à 15~h~30, heure de New York, la CFTC publie qui détenait quoi dans chaque marché à terme le mardi précédent~: producteurs et négociants, swap dealers, gestion pour compte de tiers (managed money).
Le rapport est public, bon marché et déjà ancien au moment où il paraît, et il est étudié depuis aussi longtemps qu'il existe.
De Roon, Nijman et Veld ont montré que les positions des hedgers aident à expliquer les rendements des contrats à terme sur vingt marchés~: les spéculateurs qui prennent
la contrepartie de la demande de couverture sont rémunérés pour cela. Sur les matières premières synthétiques de ce chapitre, où cette rémunération est intégrée au modèle, la position des hedgers a une IC de rang de 0.044
avec les rendements de la semaine suivante lorsqu'elle est connue immédiatement, de 0.038 lorsqu'elle est connue le vendredi, et de 0.017 trois semaines plus tard. Sur neuf ans de positionnement réel sur le maïs, le signal le plus étudié du marché
n'est pas significatif. La construction correspondante est \texttt{firm.posisig}.

\section{Rapports de positionnement}

\begin{dated}{2026-09}{Le calendrier des Commitments of Traders}\label{dat:s1:positioning-and-sentiment:cot}
Le rapport Commitments of Traders de la CFTC est en général publié chaque vendredi à 15~h~30, heure de l'Est, avec les positions arrêtées au
mardi qui précède immédiatement~; la Commission indique qu'il lui faut trois jours pour traiter les données, et les jours fériés décalent le calendrier, qu'elle publie à l'avance. Le livre~3,
chapitre~9 décrit les catégories de négociants du rapport.
\end{dated}

\begin{definition}[Pression de couverture, signal de positionnement]\label{def:s1:positioning-and-sentiment:hedging}
La \emph{pression de couverture}\index{pression de couverture} est la position nette des hedgers sur un marché à terme, généralement nette vendeuse pour les producteurs, rapportée
aux positions ouvertes~; les spéculateurs qui prennent la contrepartie perçoivent une prime de risque d'autant plus élevée que cette pression est forte. Un \emph{signal de positionnement}\index{signal de positionnement} est un
signal de trading construit à partir des positions déclarées d'une catégorie de négociants, telles qu'elles sont connues au moment de la transaction.
\end{definition}

Ce constat est plus ancien que les autres signaux du chapitre. Bessembinder a montré que les rendements des contrats à terme sur devises et sur produits agricoles varient avec les positions nettes des hedgers
une fois contrôlé le risque systématique~; de Roon et ses coauteurs ont montré que la \omterm{def:s1:positioning-and-sentiment:hedging}{pression de couverture} propre à un marché, et celle des marchés voisins de son groupe,
influent sur les rendements, et que cet effet subsiste lorsqu'on contrôle la \omterm{def:s1:flows:fit}{pression sur les prix}. La logique est celle de la vision assurantielle des contrats à terme~: les hedgers paient pour
se décharger de leur risque, et celui qui le porte est rémunéré.

Dans les positions réelles, les deux côtés sont des images miroir. Dans les données désagrégées de la CFTC pour le maïs du CBOT, de 2016 à 2026, la position nette de la gestion pour compte de tiers a varié
de $-22.8~\%$ à $+24.4~\%$ des positions ouvertes, et sa corrélation avec la position nette des producteurs et négociants est de $-0.97$
(\cref{fig:s1:positioning-and-sentiment:corn})~: les spéculateurs absorbent le risque que les hedgers leur transfèrent.

\begin{omfigure}
\begin{tikzpicture}
\begin{axis}[width=11cm,height=5.2cm,xlabel={\small année},ylabel={\small position nette (\% des positions ouvertes)},tick label style={font=\footnotesize},
  xmin=2016,xmax=2026.8,ymin=-55,ymax=30,grid=major,grid style={gray!25},legend columns=2,legend style={font=\scriptsize,at={(0.5,-0.3)},
  anchor=north,draw=none,fill=none,/tikz/every even column/.append style={column sep=8pt}},table/col sep=comma,/pgf/number format/1000 sep={}]
\addplot[omThm,thick] table[x=year,y=mm]{figdata/strategies-1/24-positioning-and-sentiment/corn.csv};
\addplot[omDef,thick] table[x=year,y=pm]{figdata/strategies-1/24-positioning-and-sentiment/corn.csv};
\addplot[black!40,dashed] coordinates {(2016,0) (2026.8,0)};
\legend{gestion pour compte de tiers,producteurs et négociants}
\end{axis}
\end{tikzpicture}

\omcaption{Contrats à terme sur maïs du CBOT~: positions nettes de la gestion pour compte de tiers et des producteurs, négociants, transformateurs et utilisateurs, en part des positions ouvertes, hebdomadaires
(mardi), de janvier 2016 à septembre 2026, d'après les Commitments of Traders désagrégés de la CFTC. Données~: \texttt{s1\_posisig.corn}.}\label{fig:s1:positioning-and-sentiment:corn}
\end{omfigure}

La position de la gestion pour compte de tiers sur le maïs permet-elle de prédire le prix du maïs~? Le test utilise le z-score sur 52 semaines de la position, tel que publié à la fin de chaque mois, et la variation du prix du maïs IMF
sur le mois qui suit le suivant (moyennes mensuelles~; sauter un mois évite que le mois propre au signal n'entre dans le rendement). Sur 113 mois,
de 2017 à 2026, la corrélation vaut 0.098 ($t = 1.04$). Après des positions acheteuses record (z-score supérieur à 1), la variation de prix des trois mois suivants a été en moyenne
de $+4.6~\%$ (28 cas), après des positions vendeuses record de $+0.2~\%$ (24), et de $-0.9~\%$ sinon (59)~: tout au plus les positions extrêmes des spéculateurs ont-elles été suivies de mouvements
dans leur sens, comme le voudrait la thèse de la \omterm{def:s1:positioning-and-sentiment:hedging}{pression de couverture}, mais avec cent mois de données, les éléments ne permettent de trancher dans aucun sens.

\section{La valeur d'un rapport, et son ancienneté}

\texttt{firm.posisig} intègre le mécanisme aux dix matières premières de \texttt{firm.synthfut} (\cref{lst:s1:positioning-and-sentiment:model}). La position nette
vendeuse des hedgers suit un AR(1) de demi-vie 20 jours de bourse, en écarts-types~; elle ajoute une prime de 0.3 fois la volatilité du marché par an et par écart-type
au rendement du lendemain. La position acheteuse nette déclarée des spéculateurs est la \omterm{def:s1:positioning-and-sentiment:hedging}{pression de couverture}, plus un demi-écart-type par unité de leur tendance sur 60 jours,
plus du bruit. Les positions sont mesurées le cinquième jour de chaque semaine (le «~mardi~») et connues après un délai de publication.
Le signal est mesuré chaque jour par rapport aux rendements des cinq jours suivants, sur les dix marchés.

\begin{center}
\small
\begin{tabular}{@{}lcccc@{}}
\toprule
délai de publication (jours de bourse) & 0 & 3 (vendredi) & 10 & 21\\
\midrule
position des hedgers~: IC de rang & 0.044 & 0.038 & 0.029 & 0.017\\
\quad statistique t & 5.1 & 4.4 & 3.4 & 2.0\\
position des spéculateurs~: IC de rang & 0.043 & 0.039 & 0.031 & 0.024\\
\bottomrule
\end{tabular}
\end{center}

Le décalage du mardi au vendredi fait perdre 14~\% du signal (\cref{fig:s1:positioning-and-sentiment:lags})~; trois semaines en font perdre 62~\%. Les deux résultats découlent de la
persistance de la pression~: avec une demi-vie de 20 jours, trois jours de vieillissement conservent 90~\% de la pression, et 21 jours en conservent 48~\%. La position des spéculateurs perd sa valeur plus
lentement, car une partie d'entre elle relève de la poursuite de tendance, et les tendances intégrées à \texttt{firm.synthfut} sont lentes~: la composante tendancielle du signal vieillit bien
même si elle n'a rien à voir avec la couverture.

\begin{omfigure}
\begin{tikzpicture}
\begin{axis}[width=10cm,height=5cm,xlabel={\small délai de publication (jours de bourse)},ylabel={\small IC de rang, cinq jours suivants},
  tick label style={font=\footnotesize},xmin=-1,xmax=22,ymin=0,ymax=0.05,grid=major,grid style={gray!25},legend pos=north east,
  legend style={font=\scriptsize},scaled y ticks=false,yticklabel style={/pgf/number format/fixed,/pgf/number format/precision=2},
  table/col sep=comma]
\addplot[omThm,thick,mark=*] table[x=lag,y=hedge]{figdata/strategies-1/24-positioning-and-sentiment/lags.csv};
\addplot[omDef,thick,mark=square*] table[x=lag,y=spec]{figdata/strategies-1/24-positioning-and-sentiment/lags.csv};
\legend{vente nette des hedgers,achat net des spéculateurs}
\end{axis}
\end{tikzpicture}

\omcaption{Signaux de positionnement sur les matières premières synthétiques à \omterm{def:s1:positioning-and-sentiment:hedging}{pression de couverture} intégrée (demi-vie 20 jours)~: l'IC de rang avec les rendements des cinq jours suivants
des positions déclarées des hedgers et des spéculateurs, selon le délai entre la mesure et sa publication. Données~:
\texttt{s1\_posisig.lag\_table}.}\label{fig:s1:positioning-and-sentiment:lags}
\end{omfigure}

La règle contrarienne, qui consiste à vendre lorsque les spéculateurs sont à un record acheteur, échoue dans le modèle par construction~: après un z-score sur 52 semaines supérieur à 1.5, le rendement du trimestre suivant
a été en moyenne de $+0.045$ volatilités annuelles, après un z-score inférieur à $-1.5$ il a été en moyenne de $-0.093$, et de $-0.034$ sinon. Une règle contrarienne suppose un
mécanisme dans lequel ce sont les spéculateurs qui paient, en provoquant un encombrement (crowding) dans des positions qui se débouclent ensuite~; le modèle n'en comporte aucun, et ni les données sur le maïs ni les
sources consultées pour ce chapitre n'en fournissent un.

\section{Enquêtes et indices de sentiment}

\begin{definition}[Indice de sentiment]\label{def:s1:positioning-and-sentiment:sentiment}
Un \emph{indice de sentiment}\index{indice de sentiment} résume l'optimisme ou le pessimisme des investisseurs à partir d'enquêtes ou de substituts fondés sur le marché (émissions, flux vers les fonds,
rendements du premier jour des nouvelles émissions, écarts de valorisation), généralement sous la forme de la première composante principale de plusieurs substituts.
\end{definition}

Baker et Wurgler ont montré ce que le sentiment prédit dans la coupe transversale des actions~: lorsque le sentiment est bas au début d'une période, les actions les plus difficiles à valoriser
et à arbitrer (petites, jeunes, volatiles, non rentables, ne versant pas de dividende, titres de croissance extrême et titres en difficulté) rapportent ensuite relativement davantage~; lorsqu'il est élevé,
elles rapportent relativement moins. Le sentiment est davantage une variable de conditionnement qu'un signal~: il indique quelles actions une vague d'optimisme a le plus éloignées de leur valeur,
et quand leur mauvaise valorisation risque d'être importante. Il évolue lentement et est publié avec un retard, si bien que son ancienneté
compte moins que celle d'un rapport de positionnement hebdomadaire.

Le positionnement sur options constitue la troisième source~: \omterm{def:s1:options-implied-signals-for-stocks:volume}{ratios put--call}, positions ouvertes par strike et gamma probable des dealers. Le chapitre~15 a mesuré ce que la demande informée sur options
dit des actions~; le chapitre~13 a mesuré ce que fait la couverture des dealers près de la clôture. En tant que positionnement, les options ajoutent le point de vue des traders qui paient
pour la convexité, publié chaque jour et par strike.

\section{Fiches de stratégie}

\begin{strategyfile}{Signal de pression de couverture}\label{strat:s1:positioning-and-sentiment:hedging}
\sfield{Qui paie, et pourquoi} Les hedgers qui versent une prime pour se décharger du risque de prix.
\sfield{Instruments et plateformes} Contrats à terme sur matières premières et sur instruments financiers, avec rapport COT.
\sfield{Signal} Position nette vendeuse des hedgers, en part des positions ouvertes, telle que publiée (le vendredi), au sein d'un marché ou à travers un groupe.
\sfield{Dimensionnement et exécution} Acheter les marchés à forte \omterm{def:s1:positioning-and-sentiment:hedging}{pression de couverture} et vendre les marchés à faible pression~; hebdomadaire, pondéré par la volatilité.
\sfield{Coûts} Faibles~; rééquilibrage hebdomadaire.
\sfield{Comment elle s'éteint} Davantage de capitaux spéculatifs se disputent la contrepartie de la couverture~; changements dans les programmes de couverture.
\sfield{Horizon, capacité, infrastructure} Semaines~; les archives COT avec leurs dates de publication.
\sfield{Backtest en toute honnêteté} Positions prises seulement à partir du vendredi de chaque rapport (plus tard après les jours fériés), jamais à partir du mardi.
\sfield{Sources} Bessembinder (1992)~; de Roon, Nijman et Veld (2000)~; ce chapitre~: IC 0.038 tel que publié, synthétique.
\end{strategyfile}

\begin{strategyfile}{Contrarien sur positionnement extrême}\label{strat:s1:positioning-and-sentiment:extreme}
\sfield{Qui paie, et pourquoi} Seulement si les spéculateurs se massent dans des positions puis les débouclent ensemble~; sinon, personne.
\sfield{Instruments et plateformes} Contrats à terme avec rapport COT.
\sfield{Signal} Position nette de la gestion pour compte de tiers (managed money) à un extrême pluriannuel.
\sfield{Dimensionnement et exécution} Prendre le contre-pied de l'extrême, en petite taille, avec un stop.
\sfield{Coûts} Faibles.
\sfield{Comment elle s'éteint} Elle peut ne pas survivre~: dans le maïs, sur 2017--2026, les extrêmes ont été suivis de mouvements favorables aux spéculateurs.
\sfield{Horizon, capacité, infrastructure} Mois.
\sfield{Backtest en toute honnêteté} Extrêmes mesurés sur les seules données passées~; dates de publication~; assez d'extrêmes pour tester.
\sfield{Sources} Aucun résultat à l'appui n'a été vérifié~; test sur le maïs de ce chapitre ($t = 1.04$ pour le signe opposé).
\end{strategyfile}

\begin{strategyfile}{Signal de positionnement sur options}\label{strat:s1:positioning-and-sentiment:options}
\sfield{Qui paie, et pourquoi} Les traders qui paient pour se protéger ou pour l'effet de levier, dont la demande se lit dans les volumes d'options et les positions ouvertes.
\sfield{Instruments et plateformes} Actions et indices avec options cotées.
\sfield{Signal} \omterm{def:s1:options-implied-signals-for-stocks:volume}{Ratios put--call} et positions ouvertes par strike, conditionnés par le sentiment.
\sfield{Dimensionnement et exécution} Biais de position sur actions ou indices~; quotidien.
\sfield{Coûts} Modérés.
\sfield{Comment elle s'éteint} Flux de couverture qui ressemble à de la spéculation~; encombrement (crowding).
\sfield{Horizon, capacité, infrastructure} De quelques jours à quelques semaines~; données d'options par strike.
\sfield{Backtest en toute honnêteté} Positions ouvertes telles que publiées le lendemain matin~; volumes orientés seulement si les données le permettent.
\sfield{Sources} Chapitre~15~; Baker et Wurgler (2006) pour le sentiment comme variable de conditionnement.
\end{strategyfile}

\section{Tutoriel~: enregistrer une position longue}\label{tut:s1:positioning-and-sentiment:tut}

\begin{tutorial}
\textbf{Objectif.} Introduire une \omterm{def:s1:positioning-and-sentiment:hedging}{pression de couverture} dans des matières premières synthétiques, mesurer ce que vaut le \omterm{def:s1:positioning-and-sentiment:hedging}{signal de positionnement} à chaque délai de publication et aux extrêmes,
et tester la série réelle de positionnement sur le maïs. \textbf{État final~:} le tableau et les deux figures.
\begin{tutsteps}
\item \textbf{Rapports et marché}~: le calendrier de publication et la \omterm{def:s1:positioning-and-sentiment:hedging}{pression de couverture} introduite.
\omcode{code/firm/posisig/firm_posisig.py}{26}{58}{Le calendrier de publication et un marché avec pression de couverture.}\label{lst:s1:positioning-and-sentiment:model}
\item \textbf{Maïs réel}~: le z-score de la gestion pour compte de tiers, tel que publié en fin de mois, face aux prix du maïs.
\omcode{code/strategies-1/24-positioning-and-sentiment/python/s1_posisig.py}{74}{102}{Positionnement sur le maïs par rapport au prix du maïs.}\label{lst:s1:positioning-and-sentiment:corn}
\item \textbf{Exécuter} \texttt{lag\_table()}, \texttt{extremes()} et \texttt{corn()}, ainsi que \texttt{fig\_posisig.py}.
\end{tutsteps}
\textbf{Pour aller plus loin.} Introduire un mécanisme d'encombrement (les positions des spéculateurs faisant \omterm{def:s1:flows:fit}{pression sur les prix}, qui finissent par se corriger) et observer apparaître la règle contrarienne~;
raccourcir la demi-vie de la pression à cinq jours, et mesurer le coût de délai du vendredi~; ajouter les swap dealers comme troisième catégorie.
\end{tutorial}

\section{Construction~: signaux de positionnement}\label{bld:s1:positioning-and-sentiment:posisig}

\begin{build}
\bfield{Rôle} Rapports tels que connus chaque jour, un marché avec \omterm{def:s1:positioning-and-sentiment:hedging}{pression de couverture}, rendements futurs, IC de rang et z-scores.
\bfield{Interface} \texttt{published(x, every, offset, lag)}, \texttt{hedging\_market(r, vol, half\_life, premium, chase, noise, rng)},
\texttt{forward(r, h)}, \texttt{ic(signal, fwd, step, overlap)}, \texttt{zscore(x, window)}~; \texttt{firm.cot}~: son \texttt{net\_share} pour les rapports réels.
\bfield{Règles} Un rapport n'est utilisable qu'à partir de son jour de publication~; les IC dont les horizons se chevauchent ont leurs t-statistiques corrigées.
\bfield{Tests d'acceptation} \texttt{code/firm/posisig/tests/}~: un calendrier mardi-vendredi vérifié à la main, sommes des rendements futurs, IC et sa correction de chevauchement, la
prime introduite.
\bfield{Pour aller plus loin} Le calendrier de publication effectif de la CFTC, avec les décalages de jours fériés~; la \omterm{def:s1:positioning-and-sentiment:hedging}{pression de couverture} entre marchés au sein des groupes~; un mécanisme d'encombrement.
\end{build}

\begin{omsources}
\item F.~A.~de Roon, T.~E.~Nijman et C.~Veld, «~Hedging pressure effects in futures markets~», \emph{Journal of Finance} 55(3), 2000.
\item H.~Bessembinder, «~Systematic risk, hedging pressure, and risk premiums in futures markets~», \emph{Review of Financial Studies} 5(4), 1992.
\item M.~Baker et J.~Wurgler, «~Investor sentiment and the cross-section of stock returns~», \emph{Journal of Finance} 61(4), 2006.
\item US Commodity Futures Trading Commission, Commitments of Traders (calendrier de publication et rapports désagrégés)~; IMF, prix mondial du maïs (via
FRED).
\end{omsources}

\section{Exercices}

\begin{exercise}[$\star$]\label{exo:s1:positioning-and-sentiment:1}
Une position mesurée le mardi est publiée le vendredi. Déterminer le premier jour à partir duquel un backtest peut l'utiliser, et expliquer ce que change un usage dès le mardi aux résultats.
\end{exercise}

\begin{exercise}[$\star$]\label{exo:s1:positioning-and-sentiment:2}
La \omterm{def:s1:positioning-and-sentiment:hedging}{pression de couverture} suit un AR(1) d'une demi-vie de 20 jours de bourse. Calculer la part qui en subsiste après trois jours, dix jours et vingt et un jours.
\end{exercise}

\begin{exercise}[$\star$]\label{exo:s1:positioning-and-sentiment:3}
Expliquer pourquoi les positions de la gestion pour compte de tiers et celles des producteurs sur le maïs sont presque des images miroir parfaites.
\end{exercise}

\begin{exercise}[$\star\star$]\label{exo:s1:positioning-and-sentiment:4}
Avec 113 mois, calculer la taille qu'une corrélation doit atteindre pour être significative à 5~\%, et dire si celle du maïs, 0.098, l'est.
\end{exercise}

\begin{exercise}[$\star\star$]\label{exo:s1:positioning-and-sentiment:5}
Expliquer pourquoi le signal des spéculateurs conserve davantage de sa valeur que celui des hedgers aux longs délais.
\end{exercise}

\begin{exercise}[$\star\star$]\label{exo:s1:positioning-and-sentiment:6}
Expliquer pourquoi la règle contrarienne échoue dans le modèle, et dire ce qu'il faudrait introduire pour qu'elle fonctionne.
\end{exercise}

\begin{exercise}[$\star\star\star$]\label{exo:s1:positioning-and-sentiment:7}
\emph{Programmation.} Exécuter \texttt{lag\_table(5.0)}, avec une demi-vie de cinq jours au lieu de vingt. Déterminer la part de l'IC des hedgers que le délai du vendredi
coûte désormais.
\end{exercise}

\begin{exercise}[$\star\star\star$]\label{exo:s1:positioning-and-sentiment:8}
\emph{Trouver la faille.} «~Nous avons construit un \omterm{def:s1:positioning-and-sentiment:sentiment}{indice de sentiment} à partir d'enquêtes, et il a prédit le rendement du mois suivant du marché avec $t = 3$ sur 1990--2020, en utilisant
les valeurs d'enquête telles que datées.~»
\end{exercise}

\section{Problème~: Record de positions longues}

\begin{problem}[{Problème du week-end --- valeur d'un rapport et son âge}]\label{pb:s1:positioning-and-sentiment:1}
Les matières premières synthétiques du chapitre, les rapports sur le maïs et le registre public.

\medskip\noindent\textbf{Partie I --- Rapports.}
\begin{enumerate}
\item À quel moment le rapport COT est-il publié, et pour quel jour~?
\item Définir la \omterm{def:s1:positioning-and-sentiment:hedging}{pression de couverture} et un \omterm{def:s1:positioning-and-sentiment:hedging}{signal de positionnement}.
\item Que trouvent Bessembinder et de Roon, Nijman et Veld~?
\item Comment les positions de la gestion pour compte de tiers sur le maïs se rapportent-elles à celles des producteurs~?
\end{enumerate}

\medskip\noindent\textbf{Partie II --- Maïs.}
\begin{enumerate}[resume]
\item Décrire le test sur le maïs et la raison pour laquelle il saute un mois.
\item Donner sa corrélation et sa statistique de Student.
\item Qu'est-ce qui a suivi les records de positions longues et les records de positions courtes~?
\item Que permettent de conclure cent mois d'observations~?
\end{enumerate}

\medskip\noindent\textbf{Partie III --- Le modèle.}
\begin{enumerate}[resume]
\item Décrire la \omterm{def:s1:positioning-and-sentiment:hedging}{pression de couverture} introduite et la position des spéculateurs.
\item Donner l'IC des hedgers pour chaque délai de publication.
\item Expliquer le coût du délai du vendredi à partir de la demi-vie.
\item Pourquoi la règle contrarienne échoue-t-elle ici~?
\end{enumerate}

\medskip\noindent\textbf{Partie IV --- Le verdict.}
\begin{enumerate}[resume]
\item Énoncer le \emph{résultat nommé}~: l'IC du \omterm{def:s1:positioning-and-sentiment:hedging}{signal de positionnement} et l'effet du délai de publication du mardi au vendredi.
\item Définir un \omterm{def:s1:positioning-and-sentiment:sentiment}{indice de sentiment}~; que trouvent Baker et Wurgler~?
\item En quoi le sentiment s'utilise-t-il différemment du positionnement~?
\item Qu'apporte le positionnement sur options~?
\item Comment backtester honnêtement une stratégie COT~?
\item Quelle fiche de stratégie ne comporte pas de mécanisme~?
\item Quel est le lien entre ce chapitre et le carry du chapitre~20~?
\item En une phrase~: qui paie un trader de positionnement~?
\end{enumerate}
\end{problem}

\section{Questions d'entretien}

\begin{interviewq}[$\star$ \iqroles{chercheur}]\label{iq:s1:positioning-and-sentiment:1}
Que contient le rapport Commitments of Traders (COT), et que ne peut-il pas vous indiquer~?
\end{interviewq}

\begin{interviewq}[$\star\star$ \iqroles{chercheur}]\label{iq:s1:positioning-and-sentiment:2}
Pourquoi les spéculateurs pourraient-ils gagner une prime sur les marchés de contrats à terme~?
\end{interviewq}

\begin{interviewq}[$\star\star$ \iqroles{trader}]\label{iq:s1:positioning-and-sentiment:3}
La gestion pour compte de tiers détient une position longue record sur le pétrole brut. Que faites-vous~?
\end{interviewq}

\begin{interviewq}[$\star\star$ \iqroles{développeur}]\label{iq:s1:positioning-and-sentiment:4}
Comment stockeriez-vous les données COT pour que les backtests n'utilisent chaque rapport qu'à partir du moment où il était public~?
\end{interviewq}

\begin{interviewq}[$\star\star$ \iqroles{chercheur}]\label{iq:s1:positioning-and-sentiment:5}
Comment construiriez-vous un \omterm{def:s1:positioning-and-sentiment:sentiment}{indice de sentiment}, et comment le testeriez-vous~?
\end{interviewq}

\begin{interviewq}[$\star\star\star$ \iqroles{chercheur}]\label{iq:s1:positioning-and-sentiment:6}
Un signal est un AR(1) de coefficient $\phi$ par jour et prédit le rendement du lendemain avec une corrélation $\rho$. Dérivez sa corrélation avec le rendement
$k$ jours après sa mesure, et la corrélation moyenne d'un rapport utilisé pendant les cinq jours suivant un délai de $\ell$ jours.
\end{interviewq}
